\section{Fragment Shader}

\subsection{Overview}
\begin{frame}{Fragment Shader Stage}
  \begin{columns}
    \small
    \begin{column}{0.7\textwidth}
      \begin{raybox}{Fragment Shader}
        \textbf{Input:} Fragments with interpolated attributes \\
        \textbf{Output:} Final fragment colors (and depth)

        \vspace{0.3cm}
        \textbf{Primary responsibility:}
        \begin{itemize}
          \item<1-> Compute final color for each fragment
          \item<2-> Apply lighting calculations
          \item<3-> Sample textures
          \item<4-> Implement material properties
          \item<5-> Output color and optional data
        \end{itemize}
      \end{raybox}
    \end{column}
    \begin{column}{0.3\textwidth}
      \begin{center}

        \begin{tikzpicture}[scale=0.7]
          % Input fragments
          \only<1->{
            \foreach \i in {0,1,2} {
              \foreach \j in {0,1,2} {
                \fill[AccentColor!30] (\i*0.4, \j*0.4 + 3) circle (0.1);
              }
            }
            \node[below] at (0.4, 2.5) {\scriptsize Fragments};
          }

          % Fragment shader box
          \only<2->{
            \node[rectangle, draw, fill=PrimaryColor!30, minimum width=1.5cm, minimum height=1cm] (fs) at (0.4,0.5) {
              \begin{minipage}{1.3cm}
                \centering
                \textbf{Fragment} \\
                \textbf{Shader}
              \end{minipage}
            };
          }

          % Output colored fragments
          \only<5->{

            \begin{scope}[yshift=-2cm]
              \foreach \i in {0,1,2} {
                \foreach \j in {0,1,2} {
                  \fill[ObjectColor] (\i*0.4, \j*0.4) circle (0.1);
                }
              }
              \node[below] at (0.4, -0.3) {\scriptsize Colored};
              \node[below] at (0.4, -0.7) {\scriptsize Fragments};
            \end{scope}
          }

          % Arrows
          \only<2->{\draw[->, thick, AccentColor] (0.4,1.8) -- (0.4,1.2);}
          \only<5->{\draw[->, thick, AccentColor] (0.4,-0.2) -- (0.4,-0.8);}
        \end{tikzpicture}
      \end{center}
    \end{column}
  \end{columns}
\end{frame}

\subsection{Input Data}
\begin{frame}{Fragment Shader Inputs}
  \begin{columns}

    \small
    \begin{column}{0.7\textwidth}
      \begin{raybox}{Interpolated Attributes}
        \textbf{From vertex shader via rasterization:}
        \begin{itemize}
          \item \textbf{Position} - World space coordinates
          \item \textbf{Normal vectors} - For lighting calculations
          \item \textbf{Texture coordinates} - UV mapping
          \item \textbf{Color} - Vertex colors
          \item \textbf{Custom attributes} - Application-specific data
        \end{itemize}
      \end{raybox}
    \end{column}
    \begin{column}{0.3\textwidth}
      \begin{tikzpicture}[scale=0.7]
        % Triangle with interpolated values
        \coordinate (A) at (0,0);
        \coordinate (B) at (3,0);
        \coordinate (C) at (1.5,2.5);

        % Triangle
        \draw[thick, ObjectColor] (A) -- (B) -- (C) -- cycle;

        % Vertex attributes
        \fill[red] (A) circle (0.1);
        \fill[green] (B) circle (0.1);
        \fill[blue] (C) circle (0.1);

        % Sample fragment
        \fill[purple] (1.2,0.8) circle (0.08);
        \node[right] at (1.3,0.8) {\tiny Fragment};

        % Interpolation lines
        \draw[dashed, gray] (A) -- (1.2,0.8);
        \draw[dashed, gray] (B) -- (1.2,0.8);
        \draw[dashed, gray] (C) -- (1.2,0.8);

        \node[below] at (A) {\tiny Red};
        \node[below] at (B) {\tiny Green};
        \node[above] at (C) {\tiny Blue};

        \node[below] at (1.5,-0.5) {\tiny Barycentric interpolation};
      \end{tikzpicture}
    \end{column}
  \end{columns}
\end{frame}

\subsection{Lighting Computation}
\begin{frame}{Per-Fragment Lighting}
  \begin{columns}
    \small
    \begin{column}{0.7\textwidth}
      \begin{raybox}{Phong Shading}
        \textbf{Phong Shading (Per-Fragment):}
        \begin{itemize}
          \item Interpolate normals, not colors
          \item Compute lighting at each fragment
          \item Higher quality, smooth specular highlights
          \item More computationally expensive
        \end{itemize}
      \end{raybox}

      \footnotesize
      \begin{mathbox}{Gouraud Shading}
        \textbf{Gouraud Shading (Per-Vertex):}
        \begin{itemize}
          \item Compute lighting at vertices
          \item Interpolate final colors
          \item Faster but lower quality
          \item Poor specular highlights
        \end{itemize}
      \end{mathbox}
    \end{column}
    \begin{column}{0.3\textwidth}
      \begin{tikzpicture}[scale=0.6]
        % Two sphere comparisons
        \begin{scope}[yshift=2cm]
          \node[circle, draw, fill=ObjectColor!30, minimum size=2cm] (sphere1) at (0,0) {};
          \node[below] at (0,-1.6) {\scriptsize Phong Shading};

          % Highlight
          \fill[white, opacity=0.8] (-0.3,0.3) ellipse (0.3 and 0.2);
        \end{scope}

        \begin{scope}[yshift=-4cm]
          \node[circle, draw, fill=ObjectColor!30, minimum size=2cm] (sphere2) at (0,0) {};
          \node[below] at (0,-1.6) {\scriptsize Gouraud Shading};

          % Faceted highlight
          \fill[white, opacity=0.6] (-0.3,0.3) -- (-0.1,0.5) -- (0.1,0.3) -- (-0.1,0.1) -- cycle;
        \end{scope}

        % Light source
        \node[circle, fill=LightColor, minimum size=0.3cm] (light) at (-3,0) {\tiny \faIcon{lightbulb}};
      \end{tikzpicture}
    \end{column}
  \end{columns}
\end{frame}

\begin{frame}[fragile]{Fragment Shader - Lighting Example}
  \begin{columns}
    \begin{column}{0.5\textwidth}
      \begin{minted}[fontsize=\scriptsize, bgcolor=LightGray!20]{glsl}
#version 330 core
// Fragment shader inputs
in vec3 FragPos;    // World position
in vec3 Normal;     // Interpolated normal
in vec2 TexCoord;   // Texture coordinates
// Outputs
out vec4 FragColor;
// Uniforms (lighting parameters)
uniform vec3 lightPos;
uniform vec3 lightColor;
uniform vec3 viewPos;
uniform sampler2D diffuseTexture;
void main() {
    // Sample texture
    vec3 albedo = texture(diffuseTexture,
                         TexCoord).rgb;
    // Normalize interpolated normal
    vec3 norm = normalize(Normal);
    // Light direction
    vec3 lightDir = normalize(lightPos
                    - FragPos);
    // Ambient
    vec3 ambient = 0.15 * albedo;
    // Diffuse (Lambert)
    float diff = max(dot(norm, lightDir),
                     0.0);
      \end{minted}
    \end{column}
    \begin{column}{0.5\textwidth}
         \begin{minted}[fontsize=\scriptsize, bgcolor=LightGray!20]{glsl}
    vec3 diffuse = diff * lightColor
                        * albedo;
    // Specular (Blinn-Phong)
    vec3 viewDir = normalize(viewPos
                    - FragPos);
    vec3 halfwayDir = normalize(
                        lightDir + viewDir);
    float spec = pow(max(
                        dot(norm, halfwayDir),
                        0.0), 64.0);
    vec3 specular = spec * lightColor;
    // Combine lighting components
    vec3 result = ambient + diffuse
                          + specular;
    FragColor = vec4(result, 1.0);
}
      \end{minted}
      \vspace{3cm}
    \end{column}
  \end{columns}
\end{frame}

\begin{frame}[fragile]{Multiple Texture Sampling Example}
      \begin{minted}[fontsize=\scriptsize, bgcolor=LightGray!20]{glsl}
// Multiple texture samplers
uniform sampler2D diffuseMap;
uniform sampler2D normalMap;
uniform sampler2D specularMap;
uniform sampler2D roughnessMap;

void main() {
    // Sample multiple textures
    vec3 albedo = texture(diffuseMap,
                         TexCoord).rgb;
    vec3 normal = texture(normalMap,
                         TexCoord).rgb;
    float specular = texture(specularMap,
                           TexCoord).r;
    float roughness = texture(roughnessMap,
                            TexCoord).r;

    // Convert normal map from [0,1] to [-1,1]
    normal = normal * 2.0 - 1.0;

    // Transform normal to world space
    // (requires TBN matrix)

    // Use in lighting calculations...
}
      \end{minted}

\end{frame}

\subsection{Output}
\begin{frame}{Fragment Shader Outputs}
  \small
  \begin{raybox}{Standard Outputs}
    \textbf{Primary output:} RGBA

    \vspace{0.3cm}
    \textbf{RGBA Color:}
    \begin{itemize}
      \item Red, Green, Blue channels (0.0 to 1.0)
      \item Alpha channel for transparency
    \end{itemize}

    \vspace{0.3cm}
    \textbf{Depth modification:}
    \begin{itemize}
      \item Can output custom depth values
      \item Used for special effects
      \item Can break early depth testing
    \end{itemize}
  \end{raybox}
\end{frame}

\subsection{Programmable Nature}
\begin{frame}{Fragment Shader - Programmable Stage}
  \small
  \begin{conceptbox}{Highly Programmable}
    \textbf{Most flexible stage} in the graphics pipeline

    \vspace{0.3cm}
    \textbf{Programmable aspects:}
    \begin{itemize}
      \item Custom lighting models (Phong, Blinn-Phong, PBR, etc.)
      \item Complex material properties and texturing
      \item Advanced effects (normal mapping, parallax mapping)
      \item Post-processing effects
      \item Procedural generation
    \end{itemize}
  \end{conceptbox}
  \pause
  \begin{mathbox}{Performance Considerations}
    \textbf{Most computationally expensive stage}
    \begin{itemize}
      \item Executes once per visible fragment
      \item Can be millions of invocations per frame
    \end{itemize}
  \end{mathbox}
\end{frame}
